\section{Decision Windows}

\begin{definition}[Decision window]
A decision window is a short interval during which the current state has lost
part of its stabilizing force while the next state has not yet been rejected.
\end{definition}

\begin{discussion}
A window exists because at least one support of the old regime has weakened
before the new regime has been ruled out. This is why windows are strategic
assets rather than motivational slogans. They temporarily compress switching
cost, reduce cue density for the current frame, or narrow the ambiguity of the
next admissible move.
\end{discussion}

\begin{proposition}[Window quality]
Not all windows are equally valuable. The highest-value windows combine a
meaningful drop in current-state stability with a low-ambiguity first action
toward the target frame.
\end{proposition}

\begin{discussion}
This criterion separates window detection from window exploitation. A person
may correctly notice that a transition moment has appeared, yet still waste it
if the first action remains underspecified. A cheap window with an unclear
entry step is often less useful than a shorter window whose first move is
already prepared.
\end{discussion}

Common windows include four operational families:
\begin{itemize}[nosep]
  \item \textbf{Temporal reset windows}: waking, post-meal reset, or the first
        minute after finishing a block. Their leverage comes from temporary cue
        sparsity before the next default routine reloads.
  \item \textbf{Spatial transition windows}: arriving at a library, gym,
        classroom, or office. Their leverage comes from entering a place where
        old cues are weaker and the new frame is easier to name.
  \item \textbf{Socially induced windows}: meetings, accountability calls, or
        public commitments. Their leverage comes from external coordination
        that makes delay more expensive and the next role more legible.
  \item \textbf{Affective or informational update windows}: sharp feedback,
        deadline shocks, or a newly perceived opportunity. Their leverage
        comes from a temporary destabilization of the old valuation structure.
\end{itemize}

\begin{example}[A weak window]
``Sit at the desk after breakfast'' often fails as a study trigger when the
document, problem set, or first visible action is still undefined. The window
is real, but ambiguity consumes it before the new state stabilizes.
\end{example}

\begin{example}[A strong window]
``Leave the office and enter the gym with clothes already packed'' often works
better because the old cues are fading, the destination already names the next
frame, and the first action requires little interpretation.
\end{example}

\begin{discussion}
The same feature that makes a window cheap also makes it perishable. If action
does not begin while the barrier is still low, the system quickly becomes
expensive again. The next section studies that expiration problem directly.
\end{discussion}
